\documentclass[10pt,a4paper]{amsart}
\usepackage[margin=19mm]{geometry}
\usepackage{amsmath,amssymb,mathtools}
\usepackage[scr=rsfs,cal=euler]{mathalfa}
\usepackage{xfrac}
\usepackage[unicode,hidelinks,bookmarks=false]{hyperref}
\newcommand{\ie}{i.e.\ }
\newcommand{\cf}{cf.\ }
\newcommand{\resp}{respectively\ }
\newcommand{\Subsection}{\subsection}
\providecommand{\nobreakdash}{\nobreak\hbox{-}}
\setlength{\emergencystretch}{2em}
\multlinegap0pt
\usepackage{mathtools}
\usepackage{amssymb}
\usepackage{tikz}
\newtheorem{theorem}{Theorem}
\newtheorem{lemma}{Lemma}
\newcommand*{\Ge}{\geqslant}
\newcommand*{\Le}{\leqslant}
\title{Sum of self-commutators of commuting operators}
\author{Sameer Chavan, Md. Ramiz Reza, Shanola S. Sequeira}
\date{Source: arXiv:2609.19287v1 (CC BY 4.0)}
\begin{document}
\maketitle

\noindent\emph{Source and licence.} Sum of self-commutators of commuting operators. Version arXiv:2609.19287v1 (CC BY 4.0), distributed by arXiv under the Creative Commons Attribution 4.0 International licence, http://creativecommons.org/licenses/by/4.0/, as recorded in the arXiv licence metadata for this exact version and shown on the abstract page https://arxiv.org/abs/2609.19287v1. Full licence notice: \texttt{licenses/arxiv-2609.19287-CC-BY-4.0.txt}.\par\smallskip

\noindent\textit{Editorial scope.} Counted unit: the article's theorem main-3 (source0.tex lines 240--248). The article's complete original proof of it is delivered with it (source0.tex lines 675--683, one \texttt{proof} environment of 9 lines, which the article places away from the statement). No mathematics is written, repaired or re-derived: the statement and the proof are the article's own lines, in the article's own notation, and no retained line is substituted or re-flowed.  Retained uncounted as the source-local context the proof needs: lines 106--113; lines 653--656.  Nothing was omitted from any retained span, so the package carries no omission record.  No retained line contains a citation, so the wrapper carries no bibliography and no bibliography entry.  No retained line contains a figure environment, an image-inclusion command or drawing code, so no image is delivered and the package declares no asset dependency.  Editorial formatting only, in addition: an AMS \texttt{amsart} preamble that restates the article's own declarations and macros used by the retained lines, namely the environments \texttt{theorem}, \texttt{lemma} on separate counters, together with the standard packages those lines need.  The retained environments print in this document as Theorem 1, Lemma 1 in order of appearance, whereas the article prints the same items with the numbering of its own full document.  The complete native source of the article as distributed in the arXiv e-print for this exact version is bundled unchanged as \texttt{sources/210766/source0.tex}.
\begin{align}
\label{d-comm-formula}
 \big[T_j, [T^*_k, T_k]\big] = \big[[T_j, T^*_k], T_k\big],
\quad 1 \Le j, k \Le d, \\
\label{d-comm-formula-im}
 \big[\Im (T_j), \Im (T_k)\big] =\frac{i}{2}\Im\big([T^*_j, T_k]\big), 
\quad 1 \Le j, k \Le d.
\end{align}

\begin{lemma} \label{gen-fact}
Let ${\bf T}$ be a commuting $d$-tuple on a complex Hilbert space $\mathcal H$ and 
let $C$ be a self-adjoint operator in $\mathcal B(\mathcal H)$ such that $CT_j=T_jC,$ $j=1, \ldots, d$, and $[{\bf T}, {\bf T}^*] \Le C.$ Then $C \Ge 0$.  
\end{lemma}

\begin{theorem} \label{main-3}
Let ${\bf T}=(T_1, \ldots, T_d)$ be a sum-hyponormal $d$-tuple on a complex Hilbert space $\mathcal H$ satisfying
\begin{equation}
\label{double-commutator}
\big[T_j, [T^*_k, T_k]\big] =0,  
\quad 1 \Le j \neq k \Le d.
\end{equation}
Then $T_1, \ldots, T_d$ are hyponormal operators. In particular, every doubly commuting sum-hyponormal $d$-tuple on a complex Hilbert space is hyponormal.
\end{theorem}

\begin{proof}[Proof of Theorem~\ref{main-3}] 	 
Consider the $(d-1)$-tuple $\widehat{\bf T}_j$ obtained from ${\bf T}$ by deleting $T_j$, and let $C=[T^*_j, T_j],$ $j=1, \ldots, d.$ By \eqref{double-commutator}, we have $CT_k = T_k C$ for every $1 \Le k \neq j \Le d$. Hence, by Lemma~\ref{gen-fact}, each of the operators $T_1, \ldots, T_d$ is hyponormal. To see the remaining half, note that by \eqref{d-comm-formula},
\begin{equation} \label{application-d-comm}
 [T_j, T^*_k]=0 \Longrightarrow \big[T_j, [T^*_k, T_k]\big] = 0
\quad 1 \Le j \neq k \Le d. 
\end{equation}
Finally, a doubly commuting $d$-tuple consisting of hyponormal operators is hyponormal, since $[\![{\bf T}^*, {\bf T}]\!]$ is the diagonal matrix whose diagonal entries are $[T^*_j, T_j],$ $j=1, \ldots, d$. This, together with \eqref{application-d-comm}, 
proves the ``in particular” statement.
\end{proof}

\end{document}
